\chapter{Stonean Spaces and Maximal Commutative Subalgebras}

Recall that a compact Hausdorff space is \textbf{Stonean} if the closure of every open set is open.
A projection in $C(K)$ is the characteristic function of a clopen subset of $K$. Stonean compact
spaces are totally separated, so locally constant functions provide finite-spectrum approximants.

\begin{lemma}
  \label{lem:fin_lin_approx_of_stonean_Sak_1_3_1} (Sak  1.3.1)
  \lean{instIsRealRankZeroContinuousMapOfTotallySeparatedSpaceOfCompactSpace,CStarAlgebra.IsRealRankZero.topologicalClosure_hull_isStarProjection}
  \leanok
  Let $K$ be a Stonean space. Then every positive self-adjoint element $a$ in $C(K)$ can be uniformly approximated by 
  finite linear combinations of projections in $C(K)$ having nonnegative coefficients. 
\end{lemma}

\begin{proof}
  \leanok
  Since a Stonean compact Hausdorff space is totally separated, locally constant self-adjoint
  functions are dense in the self-adjoint part of $C(K)$. Equivalently, $C(K)$ has real rank zero.
  In a real-rank-zero $C^{*}$--algebra, every nonnegative element lies in the closure of the
  convex cone generated by the projections. Membership in this closure is precisely uniform
  approximation by finite sums of projections with nonnegative real coefficients.
\end{proof}

In Sakai, the pointwise order on $C(K)$ in Proposition \ref{prop:stonean_of_cts_fcns_incr_cond_complete_Sak_1_3_2}
isn't specified beforehand in the book. Following $C^{*}$--algebra theory, $a\ge 0$ in $C(K)$
iff the range of $a$ is contained in $\mathbb{R}_{\ge 0}$, and $b\le a$ iff $b-a \ge 0$. The spectrum 
of a function in $C(K)$ is just its range. So $b\le a$ iff $b(t)\le a(t)$ for all $t\in K$. A net
 $(f_{\alpha})$ in $C(K)$ is \textbf{bounded} if there is a positive constant that bounds the norms of all $f_{\alpha}$.
A net is said to be \textbf{increasing} if $\alpha \le \beta$ implies $f_{\alpha}\le f_{\beta}$ for all $\alpha,\beta$
in the associated directed set. We ought to note, here, that Sakai opts not to use net language
(although he sort of does when introducing subscripts) since his directed set is identical to his net
in Proposition \ref{prop:stonean_of_cts_fcns_incr_cond_complete_Sak_1_3_2}. Recall that 
the \textbf{support} of $f\in C(K)$ is $\overline{\{x\in K | f(x)\ne 0\}}$.

\begin{lemma}
    \label{lem:cts_fns_ge_zero_le_one_directed}
    \lean{ContinuousMap.directedOn_setOf_support_subset_Icc}
    \leanok
    Let $K$ be a topological space and $U\subseteq K$. Then the set of $f\in C(K)$ supported on $U$ such that
    $0 \le f \le 1$ is a directed set w.r.t. the order on $C(K)$.
\end{lemma}

\begin{proof}
    \leanok
    The usual order is a partial order, and given $f,g\in C(K)$ supported on $U$ we have $f\vee g \in C(K)$, 
    supported on $U$, where $f\vee g(t)=\max\{f(t),g(t)\}$.
\end{proof}

\begin{proposition}
    \label{prop:stonean_of_cts_fcns_incr_cond_complete_Sak_1_3_2} (Sakai 1.3.2)
    \lean{ExtremallyDisconnected.ofConditionallyCompletePartialOrderSupContinuousMap}
    \leanok
    \uses{lem:cts_fns_ge_zero_le_one_directed}
    Let $K$ be a locally compact $R_1$ space and $\mathbb{K}$ an RCLike scalar field. Suppose every
    nonempty directed subset of $C(K,\mathbb{K})$ which is bounded above has a least upper bound.
    Then $K$ is extremally disconnected.
\end{proposition}

\begin{proof}
    \leanok
    It suffices first to treat real-valued functions; the RCLike statement follows by transporting
    directed sets through the order embedding from real-valued functions. For an open set $U$, let
    $S$ consist of the real-valued continuous functions supported in $U$ and taking values in
    $[0,1]$. Lemma \ref{lem:cts_fns_ge_zero_le_one_directed} makes $S$ directed, while $0\in S$
    and the constant function $1$ is an upper bound. Let $f$ be its least upper bound. Urysohn
    functions supported in $U$ show that $f=1$ on $U$, and hence on $\overline U$ by continuity.
    If $x\notin\overline U$, choose a continuous $g$ which is $1$ on $\overline U$ and vanishes at
    $x$. Then $fg$ is again an upper bound for $S$, so minimality gives $f\le fg$ and therefore
    $f(x)=0$. Thus $(\overline U)^c=f^{-1}(\{0\})$ is closed, proving that $\overline U$ is open.
\end{proof}

\begin{lemma}
  \label{lem:spec_masa_wstar_stonean_Sak_1_7_5} (Sak 1.7.5)
  \lean{StarSubalgebra.IsMasa.extremallyDisconnected_characterSpace}
  \leanok
  \uses{prop:stonean_of_cts_fcns_incr_cond_complete_Sak_1_3_2,def:cstar_def,thm:comm_cstar_cts_fcns_Sak_1_2_1,lem:sigma_cont_of_normal_Sak_1_7_4}
  If $C$ is any maximal commutative $C^*$--subalgebra of the $W^{*}$--algebra $M$, its
  spectrum space (maximal ideal space) is Stonean.
\end{lemma}

\begin{proof}
  \leanok
  Let $S\subseteq C$ be a nonempty directed family that is bounded above. Its image in $M$ has
  an ambient least upper bound $a$ by Lemma \ref{lem:sigma_cont_of_normal_Sak_1_7_4}. Conjugation
  by a unitary of $C$ preserves that least upper bound and fixes every member of $S$. Hence $a$
  commutes with every unitary of $C$ and therefore with all of $C$, so maximality gives $a\in C$.
  Transporting this conditional completeness through the Gelfand transform and applying
  Proposition \ref{prop:stonean_of_cts_fcns_incr_cond_complete_Sak_1_3_2} proves that the character
  space is Stonean.
\end{proof}

\begin{theorem}
  \label{thm:real_rank_zero_of_predual}
  \lean{CStarAlgebra.isRealRankZero_of_predual}
  \leanok
  \uses{lem:fin_lin_approx_of_stonean_Sak_1_3_1,lem:spec_masa_wstar_stonean_Sak_1_7_5}
  Every $C^{*}$--algebra with a specified Banach predual has real rank zero: its finite-spectrum
  self-adjoint elements are norm dense in its self-adjoint part.
\end{theorem}

\begin{proof}
  \leanok
  Given a self-adjoint element $x$, extend the commutative star subalgebra generated by $x$ to a
  masa $C$. Its character space is Stonean by Lemma
  \ref{lem:spec_masa_wstar_stonean_Sak_1_7_5}. Transporting across the Gelfand transform and using
  Lemma \ref{lem:fin_lin_approx_of_stonean_Sak_1_3_1} shows that finite-spectrum self-adjoint
  elements are dense in $C^{s}$. Their images under the inclusion $C\hookrightarrow M$ retain
  finite spectrum and approximate $x$ in the ambient norm.
\end{proof}

\begin{corollary}
  \label{cor:positive_functional_le_of_projection_le}
  \lean{PositiveLinearMap.le_on_nonneg_of_le_on_isStarProjection}
  \leanok
  \uses{def:cstar_def,lem:fin_lin_approx_of_stonean_Sak_1_3_1}
  Let $A$ be a unital real-rank-zero $C^{*}$--algebra and let $\varphi,\psi$ be positive linear
  functionals on $A$. If $\varphi(p)\le\psi(p)$ for every projection $p\in A$, then
  $\varphi(a)\le\psi(a)$ for every nonnegative $a\in A$.
\end{corollary}

\begin{proof}
  \leanok
  The inequality locus $\{a\mid\varphi(a)\le\psi(a)\}$ is a closed convex cone. It contains every
  projection by hypothesis, hence contains the cone generated by the projections. Real rank zero
  makes the closure of that cone equal to the full positive cone, which proves the claim. The
  formal theorem is more general: neither algebra need be unital, and the codomain may be any
  ordered $C^{*}$--algebra.
\end{proof}
